\documentclass{article}

%-------------------------%
% PREAMBUŁA --------------%
%-------------------------%
\usepackage[utf8]{inputenc}
\usepackage[OT4]{polski}
\usepackage{caption}
\usepackage{amsmath}
\usepackage{tikz}
\usepackage{subcaption}
\usepackage{tabularx}
\usepackage{caption}
\usepackage{array}
\usepackage{hyperref}
\usepackage{graphicx}
\usepackage{fancyhdr}
\usepackage[justification=centering]{caption}
\usepackage[margin=60pt]{geometry}
\usepackage{longtable}



%\usepackage[T1]{fontenc}
%\usepackage{amsthm}
%\theoremstyle{plain}
%\usepackage{enumitem}
%\captionsetup[table]{skip=-4pt}
%\usepackage{tabulary}
%\usepackage{sidecap}
%\usepackage{gensymb}
%\usepackage{wrapfig}
%\usepackage{lastpage}
%\usepackage{gensymb}
%\usepackage{lipsum}
%\pagestyle{fancy}


\newlength{\RoundedBoxWidth}
\newsavebox{\GrayRoundedBox}
\newenvironment{GrayBox}[1][\dimexpr\textwidth-4.5ex]%
   {\setlength{\RoundedBoxWidth}{\dimexpr#1}
    \begin{lrbox}{\GrayRoundedBox}
       \begin{minipage}{\RoundedBoxWidth}}%
   {   \end{minipage}
    \end{lrbox}
    \begin{center}
    \begin{tikzpicture}%
       \draw node[draw=black,fill=black!1,rounded corners,%
             inner sep=2ex,text width=\RoundedBoxWidth]%
             {\usebox{\GrayRoundedBox}};
    \end{tikzpicture}
    \end{center}}

%------- Typy kolumn do ustawiania szerokości ------- %
\newcolumntype{L}[1]{>{\raggedright\let\newline\\\arraybackslash\hspace{0pt}}m{#1}}
\newcolumntype{C}[1]{>{\centering\let\newline\\\arraybackslash\hspace{0pt}}m{#1}}
\newcolumntype{R}[1]{>{\raggedleft\let\newline\\\arraybackslash\hspace{0pt}}m{#1}}








\begin{document}

%-------------------------%
%Tabela nagłówkowa -------%
%-------------------------%

\begin{table}[h]
\begin{tabular}{|l|l|l|l|l|l|}
\hline
\begin{tabular}[c]{@{}l@{}}
Wydział:\\ WFiIS\end{tabular} &
\multicolumn{2}{l|}{\begin{tabular}[c]{@{}l@{}}
Imię i nazwisko:\\ 1. Axel Zuziak\end{tabular}} 
& Rok \textbf{II}         
& Grupa \textbf{B}          
& Zespół \textbf{03}      \\ \hline
\textbf{\begin{tabular}[c]{@{}l@{}}
LABOLATORIUM\\ TECHNIK \\ JĄDROWYCH\end{tabular}} &
\multicolumn{4}{l|}{ Temat:\textbf{Dozymetria promieniowania jonizującego.}
	 }                                                                                                                       & \multicolumn{1}{c|}{\begin{tabular}[c]{@{}c@{}}
Nr ćwiczenia\\ \textbf{7}
\end{tabular}} \\ \hline
\begin{tabular}[c]{@{}c@{}}
Data wykonania:\\ 06.05.2015
\end{tabular} &
\begin{tabular}[c]{@{}c@{}}
 Data oddania:\\ 20.05.2015
\end{tabular} &
\begin{tabular}[c]{@{}c@{}}
Zwrot do poprawy: \\ 
\end{tabular}                                 
& 
\begin{tabular}[c]{@{}c@{}}
Data oddania: \\ 
\end{tabular}&
\begin{tabular}[c]{@{}c@{}}
Data zaliczenia:\\
\end{tabular}& 
\begin{tabular}[c]{@{}c@{}}
OCENA:\\
\end{tabular}       
\\ & & & & & \\ \hline
\end{tabular}
\end{table}
%-----Koniec tabeli------

%%%%%%%%%%%% Wstęp %%%%%%%%%%%%
\section{Wstęp}
Celem ćwiczenia było zapoznanie się z obowiązującymi jednostkami układu SI stosowanymi w dozymetrii oraz poznanie urządzeń dozymetrycznych. W ochronie radiologicznej oraz dozymetrii stosuje się następujące wielkości:
\begin{itemize}
	\item \textbf{Dawka pochłonięta} to średnia energia tracona przez promieniowanie, pochłaniana przez ośrodek, przypadająca na jednostkę masy ośrodka. Historyczną, nieużywaną już jednostką jest rad (1rd). Obecnie używa się jednostki grej (1Gy = 1J/kg).
	\item \textbf{Dawka równoważna} jest to dawka pochłonięta w narządzie wyznaczona z uwzględnieniem rodzaju i energii promieniowania. Jednostką jest siwert.
	\item \textbf{Dawka skuteczna (efektywna)} jest to suma dawek równoważnych wyznaczonych z uzwzględnieniem czynników wagowych narządów lub tkanek.
	\item \textbf{Moc dawki pochłoniętej} jest to przyrost dawki w jednostkowym przedziale czasu. Moc dawki jest wielkością opisującą szybkość promieniowania. [Gy/s].
	\item \textbf{Dawka ekspozycyjna} promieniowania gamma lub $X$ zdefiniowana jest jako ładunek jonów jednego znaku wytwarzanych w jednostkowej masie powietrza. [C/kg]
\end{itemize}

Moc dawki $\dot{X}$ dla źródła punktowego można obliczyć ze wzoru:
\begin{equation}
	\dot{X} = \Gamma\cdot \frac{A}{l^2} \;\;\left[\frac{\text{R}}{\text{h}}\right]
	\label{moc_rad}
\end{equation}

gdzie $\Gamma$ to stała ekspozycyjna wyrażona w [Rm$^2$/Ci h].\\
Powyższą relację można zapisać następująco w układzie SI:
\begin{equation}
	\dot{X} = \Gamma_r \cdot \frac{A}{l^2} \;\;\left[\frac{\text{cGy}}{\text{h}}\right]
	\label{moc}
\end{equation}
gdzie $\Gamma_r$ jest to wartość równoważna stałe ekspozycyjnej [cGy m$^2$/s GBq]


%%%%%---------------------------- Aparatura ----------------------------%%%%%
\section{Aparatura i wykonanie ćwiczenia}

W ćwiczeniu użyto następującej aparatury:
\begin{itemize}
	\item \textbf{Radiometr Polon RKP 1-2}  - pomiary wykonywano na zakresach: 600, 200, 60 $\frac{\mu Sv}{h}$.
	\item \textbf{Mikro R meter Model 19} - pomiary wykonywano na zakresach: 5000, 500, 250 $\frac{\mu Rd}{h}$.
	\item \textbf{$^{22}Ra$} źródło o aktywności 1mCi.
	\item \textbf{Neutron Monitor NM2} licznik neutronów.
\end{itemize}
Ćwiczenie rozpoczęto od wyliczenia mocy dawki pochłoniętej zgodnie z przewidywaniami teoretycznymi \ref{moc} źródła radu w odległości 1 metra. Następnie zapoznano się z działaniem i obsługą przyrządów dozymetrycznych dostępnych w pracowni. Wykonano cechowanie dwóch powyższych przyrządów mierząc nimi moc dawki od źródła radu w osłonie 0,5 Pt. Pomiary wykonywano od 30 do 185 cm z krokiem 5 cm.


%%%%%---------------------------- Wyniki ----------------------------%%%%%
\section{Wyniki pomiarów}
\subsection{Cechowanie dozymetrów}
Zmierzone wartości dla dwóch dozymetrów naniesiono na wykresy ([\ref{polon_wykres}],[\ref{model_wykres}]). Umieszczono również wartości teoretyczne wyliczone ze wzorów \ref{moc_rad} \ref{moc}. Jako że współczynnik jakości promieniowania dla promieniowania gamma wynosi 1 nasze wyniki otrzymane w Gy/h można uznać za równoważne wynikom w Sv/h. W tabelach \ref{wyniki_polon} oraz \ref{wyniki_model} przedstawiono wyniki pomiarów. 
\begin{figure}[h!]
	\centering	
	\fontsize{8}{12}\selectfont % zmniejszam czcionke
	\resizebox{0.7\textwidth}{!}{\input{polon}}		
	\caption{Zależność równoważnika mocy dawki od odległości dla Radiometr Polon RKP 1-2.}
	\label{polon_wykres}
\end{figure}

\begin{figure}[h!]
	\centering	
	\fontsize{8}{12}\selectfont % zmniejszam czcionke
	\resizebox{0.7\textwidth}{!}{\input{model}}		
	\caption{Zależność równoważnika mocy dawki od odległości dla Mikro R meter Model 19.}
	\label{model_wykres}
\end{figure}

Następnie wykonano kalibrację urządzeń dozymetrycznych. Korzystając z funkcji fit programu $Gnuplot$ dopasowano prostą o równaniu $f(x) = ax + b$ do zależność pomiędzy przewidywaną dawką teoretyczną, a doświadczalną. 
Dla dozymetru $Mikro R$ odrzucono wyniki powyżej 2000 $\frac{\mu R}{h}$, gdyż różniły się one znacząco od przewidywań teoretycznych.

\begin{figure}[h!]
	\centering	
	\fontsize{8}{12}\selectfont % zmniejszam czcionke
	\resizebox{0.7\textwidth}{!}{% GNUPLOT: LaTeX picture with Postscript
\begingroup
  \makeatletter
  \providecommand\color[2][]{%
    \GenericError{(gnuplot) \space\space\space\@spaces}{%
      Package color not loaded in conjunction with
      terminal option `colourtext'%
    }{See the gnuplot documentation for explanation.%
    }{Either use 'blacktext' in gnuplot or load the package
      color.sty in LaTeX.}%
    \renewcommand\color[2][]{}%
  }%
  \providecommand\includegraphics[2][]{%
    \GenericError{(gnuplot) \space\space\space\@spaces}{%
      Package graphicx or graphics not loaded%
    }{See the gnuplot documentation for explanation.%
    }{The gnuplot epslatex terminal needs graphicx.sty or graphics.sty.}%
    \renewcommand\includegraphics[2][]{}%
  }%
  \providecommand\rotatebox[2]{#2}%
  \@ifundefined{ifGPcolor}{%
    \newif\ifGPcolor
    \GPcolortrue
  }{}%
  \@ifundefined{ifGPblacktext}{%
    \newif\ifGPblacktext
    \GPblacktextfalse
  }{}%
  % define a \g@addto@macro without @ in the name:
  \let\gplgaddtomacro\g@addto@macro
  % define empty templates for all commands taking text:
  \gdef\gplbacktext{}%
  \gdef\gplfronttext{}%
  \makeatother
  \ifGPblacktext
    % no textcolor at all
    \def\colorrgb#1{}%
    \def\colorgray#1{}%
  \else
    % gray or color?
    \ifGPcolor
      \def\colorrgb#1{\color[rgb]{#1}}%
      \def\colorgray#1{\color[gray]{#1}}%
      \expandafter\def\csname LTw\endcsname{\color{white}}%
      \expandafter\def\csname LTb\endcsname{\color{black}}%
      \expandafter\def\csname LTa\endcsname{\color{black}}%
      \expandafter\def\csname LT0\endcsname{\color[rgb]{1,0,0}}%
      \expandafter\def\csname LT1\endcsname{\color[rgb]{0,1,0}}%
      \expandafter\def\csname LT2\endcsname{\color[rgb]{0,0,1}}%
      \expandafter\def\csname LT3\endcsname{\color[rgb]{1,0,1}}%
      \expandafter\def\csname LT4\endcsname{\color[rgb]{0,1,1}}%
      \expandafter\def\csname LT5\endcsname{\color[rgb]{1,1,0}}%
      \expandafter\def\csname LT6\endcsname{\color[rgb]{0,0,0}}%
      \expandafter\def\csname LT7\endcsname{\color[rgb]{1,0.3,0}}%
      \expandafter\def\csname LT8\endcsname{\color[rgb]{0.5,0.5,0.5}}%
    \else
      % gray
      \def\colorrgb#1{\color{black}}%
      \def\colorgray#1{\color[gray]{#1}}%
      \expandafter\def\csname LTw\endcsname{\color{white}}%
      \expandafter\def\csname LTb\endcsname{\color{black}}%
      \expandafter\def\csname LTa\endcsname{\color{black}}%
      \expandafter\def\csname LT0\endcsname{\color{black}}%
      \expandafter\def\csname LT1\endcsname{\color{black}}%
      \expandafter\def\csname LT2\endcsname{\color{black}}%
      \expandafter\def\csname LT3\endcsname{\color{black}}%
      \expandafter\def\csname LT4\endcsname{\color{black}}%
      \expandafter\def\csname LT5\endcsname{\color{black}}%
      \expandafter\def\csname LT6\endcsname{\color{black}}%
      \expandafter\def\csname LT7\endcsname{\color{black}}%
      \expandafter\def\csname LT8\endcsname{\color{black}}%
    \fi
  \fi
  \setlength{\unitlength}{0.0500bp}%
  \begin{picture}(7370.00,5102.00)%
    \gplgaddtomacro\gplbacktext{%
      \csname LTb\endcsname%
      \put(946,704){\makebox(0,0)[r]{\strut{}$0$}}%
      \csname LTb\endcsname%
      \put(946,1393){\makebox(0,0)[r]{\strut{}$500$}}%
      \csname LTb\endcsname%
      \put(946,2082){\makebox(0,0)[r]{\strut{}$1000$}}%
      \csname LTb\endcsname%
      \put(946,2771){\makebox(0,0)[r]{\strut{}$1500$}}%
      \csname LTb\endcsname%
      \put(946,3459){\makebox(0,0)[r]{\strut{}$2000$}}%
      \csname LTb\endcsname%
      \put(946,4148){\makebox(0,0)[r]{\strut{}$2500$}}%
      \csname LTb\endcsname%
      \put(946,4837){\makebox(0,0)[r]{\strut{}$3000$}}%
      \csname LTb\endcsname%
      \put(1078,484){\makebox(0,0){\strut{}$200$}}%
      \csname LTb\endcsname%
      \put(1733,484){\makebox(0,0){\strut{}$400$}}%
      \csname LTb\endcsname%
      \put(2388,484){\makebox(0,0){\strut{}$600$}}%
      \csname LTb\endcsname%
      \put(3043,484){\makebox(0,0){\strut{}$800$}}%
      \csname LTb\endcsname%
      \put(3698,484){\makebox(0,0){\strut{}$1000$}}%
      \csname LTb\endcsname%
      \put(4353,484){\makebox(0,0){\strut{}$1200$}}%
      \csname LTb\endcsname%
      \put(5008,484){\makebox(0,0){\strut{}$1400$}}%
      \csname LTb\endcsname%
      \put(5663,484){\makebox(0,0){\strut{}$1600$}}%
      \csname LTb\endcsname%
      \put(6318,484){\makebox(0,0){\strut{}$1800$}}%
      \csname LTb\endcsname%
      \put(6973,484){\makebox(0,0){\strut{}$2000$}}%
      \put(176,2770){\rotatebox{-270}{\makebox(0,0){\strut{} Doświadczalna wartość równoważnika mocy dawki $ mu 	ext{R} / 	ext{h}$}}}%
      \put(4025,154){\makebox(0,0){\strut{}Teoretyczna wartość równoważnika mocy dawki $ mu 	ext{R} / 	ext{h}$}}%
    }%
    \gplgaddtomacro\gplfronttext{%
      \csname LTb\endcsname%
      \put(5986,4664){\makebox(0,0)[r]{\strut{}Dane pomiarowe}}%
      \csname LTb\endcsname%
      \put(5986,4444){\makebox(0,0)[r]{\strut{}f(x) = 1.38x+-133.84}}%
    }%
    \gplbacktext
    \put(0,0){\includegraphics{model_kalibracja}}%
    \gplfronttext
  \end{picture}%
\endgroup
}		
	\caption{Kalibracja dozymetru Mikro R meter Model 19.}
	\label{model_kalibracja}
\end{figure}

\begin{figure}[h!]
	\centering	
	\fontsize{8}{12}\selectfont % zmniejszam czcionke
	\resizebox{0.7\textwidth}{!}{\input{polon_kalibracja}}		
	\caption{Kalibracja dozymetru Polon RPG 1-2.}
	\label{polon_kalibracja}
\end{figure}


W powyższych równaniach współczynnik $a$ nazwany jest współczynnikiem czułości detektora, natomiast $b$ jest promieniowaniem tła. Niestety w obu przypadkach nie jest on równy jedności. 

\subsection{Wyznaczenie współczynnika osłabienia}
Umieszczono źródło promieniowania w pojemnikach ołowiowych o różnej grubości i mierzono moc równoważnika dawki w odległości 30 cm. Do punktów pomiarowych dopasowano krzywą ekspotencjalną. Współczynnik $b = 0,60$ cm$^{-1}$ jest to liniowy współczynnik osłabienia
\begin{figure}[h!]
	\centering	
	\fontsize{8}{12}\selectfont % zmniejszam czcionke
	\resizebox{0.7\textwidth}{!}{\input{oslona}}		
	\caption{Zależność równoważnika mocy dawki od grubości osłony.}
	\label{oslona}
\end{figure}
\section{Wnioski}
Obserwujemy względnie dużą zgodność wyników pomiarowych z krzywymi teoretycznymi. Moc dawki maleje wraz z kwadratem odległości od źródła. Jednak dopiero po wykonaniu kalibracji poprzez dopasowanie prostej uwidacznia się bardzo słaba dokładność szczególnie przy wysokich wartościach mocy dawki. Współczynniki czułości detektorów nie są bliskiej wartości $1$ więc nie ma podstaw do twierdzenia, że wyniki są poprawne.    

Po wykonaniu ćwiczenia nasuwa się bardzo ważny wniosek iż bardzo dobrą ochroną radiologiczną jest oddalenie się od źródła promieniowania oraz stosowanie ołowianych warstw.

\newpage
\begin{table}
\centering
\label{wyniki_polon}
\caption{Wyniki pomiarów wykonane przy pomocy dozymetru Polon RKP}
	\begin{tabular}{|C{2cm}|C{4cm}|C{2cm}|} \hline
	Odległość w [m] & Zmierzona wartość równoważnika mocy dawki [$\frac{\mu Sv}{h}$] & Wartość teoretyczna [$\frac{\mu Sv}{h}$] \\ \hline
0,3	 & 	55	 & 	87,89 \\ \hline
0,35	 & 	45	 & 	64,57 \\ \hline
0,4	 & 	37	 & 	49,44 \\ \hline
0,45	 & 	30	 & 	39,06 \\ \hline
0,5	 & 	24	 & 	31,64 \\ \hline
0,55	 & 	21	 & 	26,15 \\ \hline
0,6	 & 	19	 & 	21,97 \\ \hline
0,65	 & 	18	 & 	18,72 \\ \hline
0,7	 & 	14	 & 	16,14 \\ \hline
0,75	 & 	12	 & 	14,06 \\ \hline
0,8	 & 	11	 & 	12,36 \\ \hline
0,85	 & 	10	 & 	10,95 \\ \hline
0,9	 & 	9	 & 	9,77 \\ \hline
0,95	 & 	8	 & 	8,76 \\ \hline
1	 & 	7	 & 	7,91 \\ \hline
1,05	 & 	6	 & 	7,17 \\ \hline
1,1	 & 	5,5	 & 	6,54 \\ \hline
1,15	 & 	5	 & 	5,98 \\ \hline
1,2	 & 	4,5	 & 	5,49 \\ \hline
1,25	 & 	3,6	 & 	5,06 \\ \hline
1,3	 & 	3,4	 & 	4,68 \\ \hline
1,35	 & 	2,8	 & 	4,34 \\ \hline
1,4	 & 	2,6	 & 	4,04 \\ \hline
1,45	 & 	2,6	 & 	3,76 \\ \hline
1,5	 & 	2,4	 & 	3,52 \\ \hline
1,55	 & 	2,2	 & 	3,29 \\ \hline
1,6	 & 	2,1	 & 	3,09 \\ \hline
1,65	 & 	2,1	 & 	2,91 \\ \hline
1,7	 & 	2,2	 & 	2,74 \\ \hline
1,75	 & 	2,0	 & 	2,58 \\ \hline
1,8	 & 	2,0	 & 	2,44 \\ \hline
1,85	 & 	1,9	 & 	2,31 \\ \hline
	\end{tabular}

\end{table}


\begin{table}
\centering
\label{wyniki_model}
\caption{Wyniki pomiarów wykonane przy pomocy dozymetru Mikro R Meter.}
	\begin{tabular}{|C{2cm}|C{4cm}|C{2cm}|} \hline
	Odległość w [m] & Zmierzona wartość równoważnika mocy dawki [$\frac{\mu R}{h}$] & Wartość teoretyczna [$\frac{\mu R}{h}$] \\ \hline
0.3	 & 	4200	 & 	9166.67	\\ \hline
0.35	 & 	4100	 & 	6734.69	\\ \hline
0.4	 & 	3500	 & 	5156.25	\\ \hline
0.45	 & 	3200	 & 	4074.07	\\ \hline
0.5	 & 	2400	 & 	3300.00	\\ \hline
0.55	 & 	1900	 & 	2727.27	\\ \hline
0.6	 & 	1700	 & 	2291.67	\\ \hline
0.65	 & 	1500	 & 	1952.66	\\ \hline
0.7	 & 	1400	 & 	1683.67	\\ \hline
0.75	 & 	1250	 & 	1466.67	\\ \hline
0.8	 & 	1100	 & 	1289.06	\\ \hline
0.85	 & 	1000	 & 	1141.87	\\ \hline
0.9	 & 	900	 & 	1018.52	\\ \hline
0.95	 & 	800	 & 	914.13	\\ \hline
1	 & 	700	 & 	825.00	\\ \hline
1.05	 & 	700	 & 	748.30	\\ \hline
1.1	 & 	600	 & 	681.82	\\ \hline
1.15	 & 	500	 & 	623.82	\\ \hline
1.2	 & 	500	 & 	572.92	\\ \hline
1.25	 & 	490	 & 	528.00	\\ \hline
1.3	 & 	460	 & 	488.17	\\ \hline
1.35	 & 	420	 & 	452.67	\\ \hline
1.4	 & 	390	 & 	420.92	\\ \hline
1.45	 & 	350	 & 	392.39	\\ \hline
1.5	 & 	340	 & 	366.67	\\ \hline
1.55	 & 	320	 & 	343.39	\\ \hline
1.6	 & 	320	 & 	322.27	\\ \hline
1.65	 & 	300	 & 	303.03	\\ \hline
1.7	 & 	290	 & 	285.47	\\ \hline
1.75	 & 	270	 & 	269.39	\\ \hline
1.8	 & 	260	 & 	254.63	\\ \hline
1.85	 & 	250	 & 	241.05	\\ \hline
1.9	 & 	240	 & 	228.53	\\ \hline
		
	\end{tabular}

\end{table}


\begin{thebibliography}{9}
	
	
	\bibitem{1}
	\url{http://nucleardata.nuclear.lu.se/toi/}

	
	\bibitem{DzK}
	B. Dziunikowski, S.J. Kalita
	\emph{Ćwiczenia laboratoryjne z jądrowych metod pomiarowcyh}, Wydawnictwa AGH, Kraków 1995
	
	%\bibitem{przybycien}
	%Mariusz Przybycień \emph{Tablice Statystyczne}, dostęp on-line\\
	%\url{http://home.agh.edu.pl/~mariuszp/wfiis_stat/tablice_ps_wir.pdf}
\end{thebibliography}
\vspace{2cm}








\end{document}





%\textsc{http://nucleardata.nuclear.lu.se/toi/}
%do bibliografii






